%======================================================================
\section{¿Cómo se demuestra que una solución es buena? Teoría de la certificación}\label{sec:certteoria}
%======================================================================
Un método local entrega una solución y un certificado \emph{local} (las condiciones KKT). En un
problema no convexo eso no basta: puede haber otro mínimo mejor en otra parte. Esta sección
presenta las herramientas para obtener una garantía \emph{global}.

\subsection{Cotas superior e inferior y brecha de optimalidad}
\begin{definicion}[Brecha de optimalidad]
Para $J^\star=\min_{u\in\mathcal U}J(u)$:
\begin{itemize}
  \item una \emph{cota superior} es $UB=J(\tilde u)$ para cualquier $\tilde u\in\mathcal U$: un
        control concreto cuyo costo se conoce;
  \item una \emph{cota inferior} es un número $LB$ \emph{demostrado} tal que $J(u)\ge LB$ para
        todo $u\in\mathcal U$.
\end{itemize}
Entonces $LB\le J^\star\le UB$, y la \emph{brecha} $(UB-LB)/UB$ acota cuánto podría mejorarse
$\tilde u$.
\end{definicion}
\noindent Encontrar una cota superior es fácil: cualquier solución factible sirve. Lo difícil es
la cota inferior, porque es una afirmación sobre \emph{todos} los controles posibles, incluidos
los que nunca se evaluaron.

Hay que distinguir dos brechas que suelen confundirse:
\begin{itemize}
  \item la \textbf{brecha del solver} (residuos KKT, brecha dual de un QP) dice si un algoritmo
        resolvió bien \emph{su propio} problema;
  \item la \textbf{brecha de optimalidad} dice qué tan buena es una decisión frente a la mejor
        posible en el problema que de verdad importa.
\end{itemize}
Sólo la segunda permite comparar métodos que resuelven modelos distintos.

\begin{proposicion}[Intervalo de exceso garantizado]\label{prop:exceso}
Si $LB\le J^\star\le UB$ con $LB>0$, el exceso relativo de cualquier decisión $u$ sobre el
óptimo global cumple
\[
\frac{J(u)-UB}{UB}\ \le\ \frac{J(u)-J^\star}{J^\star}\ \le\ \frac{J(u)-LB}{LB}.
\]
\end{proposicion}
\begin{proof}
La función $s\mapsto(J(u)-s)/s=J(u)/s-1$ es decreciente en $s>0$. Basta evaluarla en $s=UB\ge
J^\star$ y en $s=LB\le J^\star$.
\end{proof}

\subsection{Relajaciones convexas}
La idea general para obtener una cota inferior es construir una función $L\le J$ en todo el
conjunto factible que sea fácil de minimizar globalmente, típicamente convexa. Su mínimo es
entonces una cota inferior de $J^\star$ (figura~\ref{fig:gap}). Las relajaciones clásicas
reemplazan cada término no convexo por una envolvente convexa~\cite{mccormick1976}.

\begin{ejemplo}[Envolvente de McCormick]
Para el término bilineal $w=xy$ con $x\in[x^L,x^U]$ e $y\in[y^L,y^U]$, desarrollando
$(x-x^L)(y-y^L)\ge0$ y $(x-x^U)(y-y^U)\ge0$ se obtiene
\[
w\ge x^Ly+y^Lx-x^Ly^L,\qquad w\ge x^Uy+y^Ux-x^Uy^U .
\]
Son dos desigualdades lineales válidas en toda la caja. Cuanto más pequeña la caja, más ajustada
la envolvente.
\end{ejemplo}
\noindent Esa última observación es la base de la \textbf{ramificación y acotamiento}
(\emph{branch and bound})~\cite{horst1996}: se divide la caja en subcajas, se calcula una cota
inferior en cada una, se descartan las subcajas cuya cota supera la mejor solución conocida, y
se subdividen las restantes hasta que la brecha es tan pequeña como se quiera. El costo es que,
en muchas dimensiones, el número de subcajas puede crecer exponencialmente.

\subsection{Aritmética de intervalos}
Para construir cotas \emph{demostradas} sobre funciones complicadas, como la salida de una
simulación, se usa la aritmética de intervalos~\cite{moore2009}. Un intervalo es
$[a]=[\underline a,\overline a]$, y las operaciones se definen para que el resultado contenga
todos los valores posibles:
\[
[a]+[b]=[\underline a+\underline b,\ \overline a+\overline b],\qquad
[a]-[b]=[\underline a-\overline b,\ \overline a-\underline b],\qquad
[a]\cdot[b]=[\min P,\ \max P],
\]
con $P=\{\underline a\underline b,\ \underline a\overline b,\ \overline a\underline b,\ \overline a\overline b\}$.

\begin{proposicion}[Propiedad de inclusión]
Si una función $F$ se evalúa con aritmética de intervalos sobre $[x]$, el intervalo resultante
contiene $\{F(x):x\in[x]\}$.
\end{proposicion}
\noindent La demostración se hace operación por operación: cada operación contiene todos sus
resultados posibles, y la composición conserva esa propiedad. Si además cada extremo se redondea
\emph{hacia afuera}, el redondeo de la computadora también queda cubierto, y el resultado es una
cota \emph{rigurosa}. El paquete \code{IntervalArithmetic.jl}~\cite{intervaljl} implementa todo
esto en Julia.

\begin{ejemplo}[El problema de dependencia]
Con $[x]=[1,2]$, la expresión $x-x$ vale exactamente cero, pero la aritmética de intervalos da
$[1,2]-[1,2]=[-1,1]$: no sabe que los dos $x$ son el mismo número. La cota es correcta, pero
pesimista. En una simulación de 26 pasos y 17 estados, esa sobreestimación se acumula paso a
paso y puede volver inútil la cota.
\end{ejemplo}
\noindent El remedio clásico es la \textbf{forma del valor medio}: en lugar de evaluar $F([x])$,
se evalúa
\[
F([x])\subseteq F(\check x)+[F']\big([x]\big)\,\big([x]-\check x\big),
\]
con $\check x$ el centro del intervalo y $[F']$ la derivada evaluada con intervalos. Para $F(x)=x-x$,
la derivada es cero y la forma da exactamente $[0,0]$. Más en general, los \textbf{modelos de
Taylor} representan cada cantidad como un polinomio con coeficientes numéricos, que conserva la
correlación entre variables, más un \emph{resto} encerrado en un intervalo~\cite{makino2003}.
Ésa es la estrategia del paso 4b (sección~\ref{sec:certificacion}): la parte lineal de la
respuesta del reactor se calcula exactamente, y sólo el resto no lineal se encierra con
intervalos.

\subsection{Verificar que una solución existe: punto fijo}
Cuando una cantidad se define implícitamente, como el estado de un paso de Euler implícito, que
aparece a ambos lados de su ecuación, hay que demostrar además que la solución existe dentro de
la caja propuesta. La herramienta es el teorema de punto fijo de Brouwer: si una función
continua lleva una caja dentro de sí misma, tiene un punto fijo en ella. El operador de
Krawczyk~\cite{krawczyk1969} convierte esa idea en una prueba computable para sistemas de
ecuaciones: si el operador, evaluado con intervalos sobre la caja, cae estrictamente dentro de
ella, la caja contiene una solución y ésta es única.

\takeaway{\textbf{Para recordar.} Una garantía global requiere una cota inferior demostrada. Se
obtiene con una relajación convexa del problema. Cuando el problema contiene una simulación, la
relajación necesita acotar con rigor cuánto se aparta el modelo no lineal de una aproximación
simple, y eso se hace con aritmética de intervalos y modelos de Taylor, verificando por punto
fijo que las cajas propuestas contienen la solución.}
